\section{Introduction}
\label{sec:intro}

Hardware enclaves (e.g. Intel SGX, Intel TDX, AMD SEV-SNP) provide strong hardware-rooted memory encryption and isolated execution environments \cite{costan2016intel}. These Trusted Execution Environments (TEEs) enable cloud customers to process sensitive workloads without trusting the underlying host operating system or cloud hypervisor.

Despite hardware isolation guarantees, TEEs are vulnerable to deterministic side-channel attacks \cite{xu2015controlled, vanbulck2018foreshadow}. An untrusted host OS can monitor page-fault sequences to reconstruct secret data processing pipelines. Furthermore, traditional remote attestation schemes expose raw binary measurement hashes ($MRENCLAVE$), allowing malicious verifiers to fingerprint specific software versions.

Existing side-channel mitigation frameworks rely on static Path ORAM schemes \cite{stefanov2013pathoram}. However, static tree traversals incur up to $100\times$ memory access overhead, rendering them impractical for production enclaves.

To resolve these limitations, we introduce \textbf{EnclaveShield}, a confidential computing framework combining Zero-Knowledge Remote Attestation with an Adaptive Dynamic ORAM Page Obfuscator.

Our primary contributions include:
\begin{enumerate}
    \item \textbf{Zero-Knowledge Remote Attestation Protocol}: Quote verification hiding raw binary measurement hashes.
    \item \textbf{Adaptive Dynamic ORAM Engine}: Frequency-aware tree balancing reducing page traversal overhead by $10.2\times$.
    \item \textbf{EnclaveBench Benchmark}: A 20-scenario benchmark evaluating 5 side-channel threat profiles across 3 random seeds.
\end{enumerate}
